% =====================================================================
\chapter{Decodificación: búsqueda y muestreo}
\label{tema:1}
% =====================================================================

Un modelo de lenguaje autorregresivo proporciona, en cada paso, una distribución $P(\cdot\mid x_{<t})$ sobre el vocabulario; la \emph{decodificación} es la regla que convierte esa sucesión de distribuciones en un texto concreto. Este tema presenta las dos familias de reglas: las de \emph{maximización} (búsqueda voraz y búsqueda en haz), que intentan encontrar la secuencia más probable, y las de \emph{muestreo} (temperatura, \emph{top-k} y núcleo), que extraen una secuencia aleatoria de una versión modificada de la distribución. Se siguen el §7.6 y el §13.4 de Jurafsky y Martin~\cite{jm3}, el §10.8 de Zhang \emph{et al.}~\cite{d2l}, la sección~3 de Holtzman \emph{et al.}~\cite{holtzman}, el §5.4 de Fan \emph{et al.}~\cite{fan} y la sección~6 de la tarea~1 de CS336~\cite{cs336a1}.

\section{El problema de la decodificación}

Sea $P$ un modelo de lenguaje sobre $\V$ con símbolo de fin \fin\ (módulo~1). Por la regla de la cadena, la probabilidad de una continuación $y=(y_1,\dots,y_T)$ de un prompt $x$ es
\[
P(y\mid x)=\prod_{t=1}^{T}P(y_t\mid x,y_{<t}),\qquad
\ln P(y\mid x)=\sum_{t=1}^{T}\ln P(y_t\mid x,y_{<t}).
\]

\begin{definicion}[Regla de decodificación]\label{def:decodificacion}
Una \emph{regla de decodificación} es un procedimiento que, a partir del prompt $x$ y del acceso a las distribuciones condicionadas $P(\cdot\mid x,y_{<t})$, produce una continuación $y$, de forma determinista o aleatoria. La \emph{decodificación de máxima probabilidad} (MAP) busca
\[
y^{\ast}=\argmax_{y}\ P(y\mid x),
\]
donde el máximo se toma sobre las secuencias terminadas en \fin.
\end{definicion}

\begin{proposicion}\label{prop:exhaustiva}
Con un vocabulario de $K$ tokens, el número de continuaciones de longitud exactamente $T$ es $K^{T}$. En consecuencia, la búsqueda exhaustiva del máximo entre las secuencias de longitud $T$ requiere evaluar $K^{T}$ secuencias.
\end{proposicion}

Con $K=5\cdot 10^{4}$ y $T=20$, $K^{T}\approx 9{,}5\cdot 10^{93}$: la decodificación MAP exacta es inabordable y se aproxima con búsquedas locales~\cite[§10.8.2]{d2l}.

\section{Búsqueda voraz y búsqueda en haz}

\begin{definicion}[Búsqueda voraz]
La \emph{búsqueda voraz} (\emph{greedy}) elige en cada paso el token más probable, $y_t=\argmax_{v\in\V}P(v\mid x,y_{<t})$, hasta emitir \fin\ o alcanzar una longitud máxima~\cite[§7.6.1]{jm3}.
\end{definicion}

\begin{proposicion}\label{prop:voraz-no-optima}
La búsqueda voraz no encuentra, en general, la secuencia de máxima probabilidad.
\end{proposicion}

\begin{definicion}[Búsqueda en haz]\label{def:haz}
Sea $m\ge 1$ la \emph{anchura del haz}. La \emph{búsqueda en haz} mantiene un conjunto $\mathcal{H}_t$ de, a lo sumo, $m$ prefijos de longitud $t$. En el paso $t+1$ forma todas las extensiones $(h,v)$ con $h\in\mathcal{H}_t$ y $v\in\V$, las puntúa con $\ln P(h,v\mid x)=\ln P(h\mid x)+\ln P(v\mid x,h)$ y conserva las $m$ de mayor puntuación. Las extensiones que terminan en \fin\ se retiran del haz como candidatas finales; al acabar se devuelve la candidata de mayor puntuación~\cite[§13.4]{jm3}, \cite[§10.8.3]{d2l}.
\end{definicion}

\addtocounter{figure}{1}%

\begin{proposicion}[Coste y casos extremos]\label{prop:coste-haz}
Para una salida de longitud $T$ y vocabulario de tamaño $K$, la búsqueda en haz de anchura $m$ evalúa a lo sumo $mKT$ extensiones. Con $m=1$ coincide con la búsqueda voraz, y con $m\ge K^{T-1}$ no poda ningún prefijo y coincide con la búsqueda exhaustiva entre las secuencias de longitud~$T$.
\end{proposicion}

\begin{observacion}[Sesgo hacia las secuencias cortas]\label{obs:longitud}
Cada término $\ln P(y_t\mid x,y_{<t})$ es menor o igual que cero, de modo que la puntuación $\ln P(y\mid x)$ no aumenta al alargar la secuencia. La búsqueda en haz tiende entonces a preferir candidatas cortas. La corrección habitual divide por una potencia de la longitud,
\[
\mathrm{punt}(y)=\frac{1}{L^{a}}\ln P(y\mid x),\qquad L=\abs{y},\ a\in[0,1],
\]
con $a=0{,}75$ en el ejemplo de~\cite[§10.8.3]{d2l}; con $a=1$ se obtiene la log-probabilidad media por token de~\cite[§13.4]{jm3}.
\end{observacion}

\begin{observacion}[Degeneración]
La búsqueda en haz es la regla habitual en traducción automática, donde la salida está muy determinada por la entrada. En la generación abierta produce textos repetitivos: Holtzman \emph{et al.} muestran que el texto humano tiene una probabilidad por token claramente menor que la del texto obtenido por maximización y que las repeticiones se refuerzan a sí mismas, porque cada repetición aumenta la probabilidad de la siguiente~\cite[§4.3]{holtzman}. Por eso los asistentes conversacionales muestrean.
\end{observacion}

\section{Muestreo con temperatura}

\begin{definicion}[Softmax con temperatura]\label{def:temperatura}
Sea $z\in\R^{K}$ el vector de logits y sea $\tau>0$. La \emph{distribución con temperatura} $\tau$ es
\[
p_\tau(i)=\frac{e^{z_i/\tau}}{\sum_{j=1}^{K}e^{z_j/\tau}},\qquad i=1,\dots,K.
\]
El \emph{muestreo con temperatura} extrae cada token de $p_\tau$~\cite[§7.6.3]{jm3}, \cite[§3.3]{holtzman}. Con $\tau=1$ se recupera la distribución del modelo.
\end{definicion}

\begin{proposicion}[Temperatura como potencia]\label{prop:potencia}
Sea $p=p_1$ la distribución del modelo. Entonces $p_\tau(i)\propto p(i)^{1/\tau}$. En particular, elevar las probabilidades a una potencia $a>0$ y renormalizar equivale a muestrear con temperatura $\tau=1/a$.
\end{proposicion}

\begin{proposicion}[Límites de la temperatura]\label{prop:limites-tau}
Sea $M=\{i:z_i=\max_j z_j\}$ el conjunto de logits máximos. Entonces
\[
\lim_{\tau\to 0^{+}}p_\tau(i)=\begin{cases}1/\abs{M}, & i\in M,\\ 0,& i\notin M,\end{cases}
\qquad
\lim_{\tau\to\infty}p_\tau(i)=\frac{1}{K}.
\]
Si el máximo es único, el límite $\tau\to 0^{+}$ es la búsqueda voraz.
\end{proposicion}

\addtocounter{definicion}{1}%

La temperatura reescala todos los logits por igual y nunca anula un token: la cola de tokens improbables conserva masa positiva para cualquier $\tau>0$. Los métodos de truncamiento la eliminan.

\section{Truncamiento: \emph{top-k} y muestreo de núcleo}

\begin{definicion}[Distribución truncada]\label{def:truncada}
Sea $p$ una distribución sobre $\V$ y sea $S\subseteq\V$ con $p(S)=\sum_{i\in S}p(i)>0$. La \emph{distribución truncada} a $S$ es $p_S(i)=p(i)/p(S)$ si $i\in S$ y $p_S(i)=0$ en otro caso.
\end{definicion}

\begin{definicion}[\emph{Top-k} y núcleo]\label{def:topk-topp}
Sea $p_{(1)}\ge p_{(2)}\ge\dots\ge p_{(K)}$ la ordenación decreciente de las probabilidades.
\begin{enumerate}[label=(\alph*)]
  \item El muestreo \emph{top-k}~\cite[§5.4]{fan} trunca $p$ al conjunto $V^{(k)}$ de los $k$ tokens más probables.
  \item El \emph{muestreo de núcleo} o \emph{top-p}~\cite[§3.1]{holtzman}, con umbral $p\in(0,1]$, trunca al menor conjunto $V^{(p)}$ tal que $\sum_{i\in V^{(p)}}p(i)\ge p$.
\end{enumerate}
\end{definicion}

\begin{proposicion}[El núcleo está formado por los más probables]\label{prop:nucleo}
Sea $k^{\ast}=\min\{k:\sum_{i=1}^{k}p_{(i)}\ge p\}$. Entonces el tamaño mínimo de un conjunto con masa al menos $p$ es $k^{\ast}$, y el conjunto de los $k^{\ast}$ tokens más probables lo alcanza. Es decir, el muestreo de núcleo coincide con \emph{top-k} con $k=k^{\ast}$, donde $k^{\ast}$ depende de la distribución de cada paso.
\end{proposicion}

\begin{proposicion}[Distorsión del truncamiento]\label{prop:tv-truncada}
Para cualquier $S$ con $p(S)>0$, la distancia de variación total entre la distribución truncada y la original es la masa eliminada:
\[
\TV(p_S,p)=\tfrac12\sum_{i\in\V}\abs{p_S(i)-p(i)}=1-p(S).
\]
\end{proposicion}

Así, el muestreo de núcleo con umbral $p$ garantiza $\TV(p_S,p)\le 1-p$ en cada paso, mientras que \emph{top-k} con $k$ fijo no controla la masa eliminada: si el modelo duda entre muchas opciones razonables, $k$ fijo descarta masa grande; si está seguro, conserva tokens de cola~\cite[§3.2]{holtzman}.

\addtocounter{definicion}{1}%

\begin{observacion}[Combinación y uso práctico]
Las reglas se componen en el orden de la \cref{fig:cadena-muestreo}: los logits se dividen por $\tau$, se aplica la softmax, se trunca con \emph{top-k} y luego con el núcleo, se renormaliza y se muestrea~\cite[§7.6.4]{jm3}. Las API públicas de Anthropic, Google y xAI exponen la temperatura y el umbral del núcleo como parámetros de la petición; los valores que usan por defecto sus aplicaciones de chat no son públicos.
\end{observacion}

\addtocounter{figure}{1}%

